\subsection{The conditional density increases at zero}

Theorem~\ref{thm:bias-continuity} needs $h_a'>0$ on the whole negative
half-line. Unimodality of $h_a$~\cite{glrs2025} and reflection
dominance give only $h_a'\ge0$ there.

\begin{lemma}[a balanced start]\label{lem:balanced-start}
For $1/2<a<1$ the density conditional on a positive first step
satisfies
\[
 h_a'(t)>0\qquad(t<0).
\]
The proof uses the walk whose first two steps are $+1,-1$. For
$0\le a\le1$ its count law at each length $n\ge2$ is log-concave, with
positive support $\{1,\ldots,n-1\}$. For $1/2<a<1$ its limiting density
$f_a^{+-}$ is symmetric and log-concave, with $(f_a^{+-})'(t)<0$ for
$t>0$.
\end{lemma}

\begin{proof}
Given $X_1=1$, the second step is $+1$ with probability $(1+a)/2$. So
$h_a$ is a mixture of the limit densities $f_a^{++}$ and $f_a^{+-}$ of
the starts $+,+$ and $+,-$. We show that $f_a^{+-}$ is strictly
increasing on the negative half-line and that $f_a^{++}$ is
nondecreasing there. We prove the log-concavity of the count law for the
start $+,-$, identify its limit, pass the log-concavity to the limit,
and then treat the start $+,+$.
We run the log-concavity induction of~\cite[Proposition~2.2]{glrs2025} at
the points where $A,B\ge d$ below. Starting
with one step in each direction keeps the support inside
$\{1,\ldots,n-1\}$, so the two ends of the support need no check. The count law given $X_1=1$ is not
log-concave~\cite{tv2026}, so this induction cannot be run from one
fixed step.
Write one update as
$g_k=Af_k+Bf_{k-1}$, where
\[
 A=\frac{1+a}{2}-\frac{ak}{n},\qquad
 B=\frac{1-a}{2}+\frac{a(k-1)}n,\qquad d=\frac an.
\]
For $u=f_{k-2}$, $v=f_{k-1}$, $w=f_k$, and $z=f_{k+1}$, expansion gives
\begin{align*}
 g_k^2-g_{k-1}g_{k+1}
 &=(A^2-d^2)(w^2-vz)+(B^2-d^2)(v^2-uw)\\
 &\quad +(A-d)(B-d)(wv-uz)+d^2(w-v)^2.
\end{align*}
For $2\le k\le n-1$ we have $A,B\ge d\ge0$.
The preceding row is log-concave with interval support, so all
three gaps are nonnegative. At the two boundaries one of the
neighboring output masses is zero, and there is nothing to check. We
start from the point mass at one and induct. The recurrence also
shows the masses are strictly positive on the support.

Next we identify the limit. Step $t+1$ copies or reverses an earlier
step $X_i$. For $t\ge1$, join step $t+1$ to $X_i$ and give the
resulting edge the sign $X_{t+1}X_i$. Doing so creates a random
recursive tree on the steps, where the edge signs are independent both
of the tree and of each other. Each of these edge signs is $+1$ with
probability $(1+a)/2$. Removing the edge from step $2$ to step $1$
splits the steps into the part of step $1$ and the part of step $2$. A
new step joins the part that contains the step it copies or reverses.
Since that step is chosen uniformly, the new step joins a part with
probability equal to the fraction of the earlier steps that lie in
that part. Hence the sizes of the two
parts form a P\'olya urn started from one ball of each color, and the
share of the part of step $1$ converges almost surely to a uniform
variable $U$. Given the sequence of parts that the steps join, each
part is a random recursive tree in its own time, with its own
independent edge signs. So the sum of the steps in the part of step
$1$ is an elephant walk with a positive first step, observed at the
size of that part. The sum in the part of step $2$ is minus an
independent copy of such a walk. Both sizes tend to infinity. By the
almost sure convergence of $S_k/k^a$ given $X_1=1$~\cite[Theorem~3.7]{bercu2018},
\[
 \frac{S_n}{n^a}\longrightarrow
 \mathcal L_a^{+-}:=U^a\mathcal L_a'-(1-U)^a\mathcal L_a''
 \qquad\text{almost surely},
\]
where $\mathcal L_a'$ and $\mathcal L_a''$ are copies of $\mathcal L_a$
and $U$, $\mathcal L_a'$, $\mathcal L_a''$ are independent. With two
positive first steps the same split gives the limit
$\mathcal L_a^{++}=U^a\mathcal L_a'+(1-U)^a\mathcal L_a''$.

These limits have entire densities. Let $\phi_a$ be the characteristic
function of $\mathcal L_a$, and take $1<r<1/a$. Since $ra<1$, we have
$v^{ra}+(1-v)^{ra}\ge1$ for $0\le v\le1$. Hence the bound of
Lemma~\ref{lem:analytic-density} gives
\[
 \bigl|\E e^{it\mathcal L_a^{+\pm}}\bigr|
 \le\int_0^1|\phi_a(tv^a)|\,|\phi_a(t(1-v)^a)|\,dv
 \le4e^{-c|t|^r}.
\]
So, exactly as in the proof of that lemma, Fourier inversion gives
densities $f_a^{+-}$ and $f_a^{++}$ that extend to entire functions.
In particular they are smooth, so Lemma~\ref{lem:derivative-limits}
applies to both starts.

Now interpolate the rescaled lattice masses log-linearly. The result
is log-concave. By Lemma~\ref{lem:derivative-limits} with $r=0$ the
rescaled masses converge locally uniformly to $f_a^{+-}$. The
log-linear interpolant lies between adjacent masses, so it has the
same limit. A pointwise limit of log-concave functions is log-concave.
Hence $f_a^{+-}$ is log-concave. Symmetry follows from the
representation of $\mathcal L_a^{+-}$. The law of $\mathcal L_a^{+-}$
is a mixture over $U$ of convolutions of two positive densities, since
$h_a>0$. So $f_a^{+-}>0$. Now $(\log f_a^{+-})'$ is nonincreasing and is zero at zero.
Suppose it were also zero at some $t>0$. Then the density would be
constant on $[0,t]$, which analyticity rules out. Hence
$(f_a^{+-})'<0$ for $t>0$.

Next take two positive initial steps. We claim the first differences
are nonnegative on the negative half-line, and use induction from the
point mass at count two. Each entry on the negative half-line has
both parents in the lower half of the previous row. The one exception is the central
difference of an odd row, and there reflection dominance gives the
sign. The induction of Proposition~\ref{prop:dominant-half} runs in
the same way from the point mass at count two. The left end causes no
trouble, because the support starts at count two and the differences
at the left end are nonnegative. The derivative limit of Lemma~\ref{lem:derivative-limits} then
gives $(f_a^{++})'(t)\ge0$ for $t<0$.
Finally
$h_a=((1+a)/2)f_a^{++}+((1-a)/2)f_a^{+-}$.
The second summand has strictly positive derivative there.
\end{proof}

Gu\'erin, Laulin, Raschel, and Simon prove positivity of the conditional
mode for memory sufficiently close to one~\cite[Remark~2.1(a)]{glrs2025}.
We need it for all $a\in(1/2,1)$, in the stronger form of a positive
slope at zero.

\begin{lemma}[the conditional mode is positive]\label{lem:positive-conditional-mode}
For every $1/2<a<1$, we have $h_a'(0)>0$. In particular the unique
conditional mode lies strictly to the right of zero.
\end{lemma}

\begin{proof}
For $z=1$ or $z=2$ the number $N^{2a}e_N(z)/z$ defined below is a
rescaled central first difference, and it tends to $4h_a'(0)$. We show
that its lower limit is positive. After division by $z$ and a product
normalization, the reflection differences become averages of earlier
rows along a backward chain. The chain stays in a diffusive band with
probability above one half, and on the band the values at a fixed row
$n_0$ are bounded below. Write $f_n(k)=\Prob_+(A_n=k)$, with $\Prob_+$ the law given $X_1=1$, and,
for $z\equiv n\pmod2$,
$e_n(z)=f_n((n+z)/2)-f_n((n-z)/2)$.
With a positive start, the reflection induction of
Proposition~\ref{prop:dominant-half} is strict for $0<z\le n$. Indeed
$e_1(1)=1$, and each noncentral entry of the upper half has a
positive parent with a positive coefficient. The recurrence is
\begin{gather*}
 e_{n+1}(z)=\alpha_n(z)e_n(z+1)+\beta_n(z)e_n(z-1),\\
 \alpha_n(z)=\frac12-\frac{a(z+1)}{2n},\qquad
 \beta_n(z)=\frac12+\frac{a(z-1)}{2n}.
\end{gather*}
Since $\alpha_n(z)(z+1)+\beta_n(z)(z-1)=z(1-2a/n)$, the function $z$
is reproduced by the recurrence up to the factor $1-2a/n$. So dividing
by $z$ and by these factors turns the recurrence into an average, and
$e_N(z)/z$ is the slope of the count law at the center.
Fix $n_0>4a$. Divide by $z$ and multiply by
$\prod_{j=n_0}^{n-1}(1-2a/j)^{-1}$. Call the new values $v_n(z)$.
They satisfy a backward averaging equation, with weights
\[
 w_+=\frac{\alpha_n(z)(z+1)}{z(1-2a/n)},\qquad
 w_-=\frac{\beta_n(z)(z-1)}{z(1-2a/n)}.
\]
In the band we use below, these weights are nonnegative and add up to
one. At $z=1$ the parent $0$ gets weight zero. Let $z_n$ be the
backward chain with these weights. Where the weights are nonnegative,
$v_n(z_n)$ is a martingale in the backward time. Then
\[
 \E[z_n^2\mid z_{n+1}=z]
 =\frac{n-4a}{n-2a}z^2+\frac{3n-4a}{n-2a}.
\]
With $\gamma_n=\Gamma(n-2a)/\Gamma(n-4a)\sim n^{2a}$, the process
$Y_n=z_n^2/\gamma_n$ has backward drift
$\delta_n=(3n-4a)/\{(n-2a)\gamma_n\}$. Note that
$\sum_{n\ge n_0}\delta_n\to0$ as $n_0\to\infty$.

Fix $R>0$. We start the chain at $z_N=1$ if $N$ is odd and at $z_N=2$
if $N$ is even. We stop it at row $n_0$, or when it leaves the band
$0<z/\sqrt{\gamma_n}\le R$. Since $\sqrt{\gamma_n}\sim n^a$ and $a<1$,
we can take $n_0$ so large that $R\sqrt{\gamma_n}+1\le n$ for all
$n\ge n_0$. Then every state in the band and its two parents lie in
the support, and $\alpha_n(z),\beta_n(z)\ge0$ there. The chain uses the
weights only at band states before it stops, so we need their
nonnegativity only there. The chain moves by $\pm1$, and it can reach
$0$ only from $z=1$, where the parent $0$ has weight zero. So it
cannot leave the band through the bottom, and on leaving the band
$Y>R^2$. Up to the stopping time, the mean of $Y$ is at most
$z_N^2/\gamma_N+\sum_{n\ge n_0}\delta_n$. Taking $n_0$ and then $N$
large, Markov's inequality shows that the chain leaves the band with
probability below one half. Every stopped value of $v$ is positive,
by the strict reflection induction. The band states of row $n_0$ are
the finitely many $z$ with $0<z\le R\sqrt{\gamma_{n_0}}$, and at row
$n_0$ the normalizing product is empty, so $v_{n_0}(z)=e_{n_0}(z)/z$.
Their minimum $m_0>0$ depends only on $n_0$ and $R$, not on $N$. The
chain stops in this row with probability above one half. The
stopping time is bounded, so $v_N(z_N)$ is the mean of the stopped
value, and $v_N(z_N)\ge m_0/2$. The
product normalizer is a positive constant times $N^{2a}$ in the
limit. Hence
\[
 \liminf_N N^{2a}\frac{e_N(z_N)}{z_N}>0.
\]
By Lemma~\ref{lem:derivative-limits} for the conditional start,
the first differences converge, and this limit is $4h_a'(0)$. For
even $N$ we use that $e_N(2)$ is the sum of two adjacent central first
differences. So
$h_a'(0)>0$. Since $h_a$ is unimodal, its mode is to the right of
zero.
\end{proof}
